%% APPENDIX D -- the false-alarm expectation and the control ensemble.
%% Owner: appendix-triage.  Carries \label{app:falsealarm}, cited by
%% 04_method.tex and 08_backmatter.tex.
\section{The false-alarm expectation and the control ensemble}
\label{app:falsealarm}

The paper's central claim is that the number of unattributed crossings is the
number chance predicts. This appendix gives the expectation that claim rests
on, what it is conditioned on, and what the control ensemble that calibrates
it can and cannot support.

\subsection{What the expectation is conditioned on}
%% (subsection label retired: nothing cites it)

The expectation is built from the control statistics themselves, which are
the only empirical null the data provide, and it imposes the condition an
event must actually satisfy. Under the disposition chain this paper applies,
that condition is to reach the trigger \emph{and} outrank every control.
Drawing one control index per execution block --- so that windows sharing
weather, calibration and correlator setup are not treated as independent
trials --- and resampling blocks, the expectation over the \NWinA{} Class~A
windows is \ClustStageMean{} unattributed rank-passing crossings against the
\ClustStageNObs{} observed, and the probability of that count or more is
$P(\geq\ClustStageNObs)=\ClustStagePobs$. Clustering changes the number
little, because the \emph{position} of a block's control maximum is not
shared across its windows.

\emph{The conditioning is restrictive, and that is what makes it
informative.} Exactly one probe per window is the ensemble maximum, so the
probability that the stellar position would have been that maximum, treated
as exchangeable with the probes, is $1/(\NCtrl+1)$ in windows whose control
maximum reaches $T_{\rm trig}=\StageNullTrig$ and zero in the rest. Only
\StageNullEligible{} of the \StageNullAllWin{} windows that retain a control
vector have a control maximum reaching the trigger at all, and
\StageNullEligibleA{} of those are Class~A. \emph{This is why every
rank-passing crossing in the survey is Class~A}, and why the choice of
reference class largely dissolves once the trigger condition is imposed. An
expectation conditioned on the rank alone predicts a different quantity ---
the number of windows in which a no-signal position would rank first --- and
is a pipeline diagnostic rather than a statement about candidates.

A second conditioning is reported beside it, on the same control ensemble,
and it is the one the body quotes: \emph{without} the rank clause ---
trigger, localisation and attribution only --- \RsevChanceTrigA{} Class~A
windows are expected to trigger by chance, the visibility fit localises
\RsevChanceFracLoc{} of events at the trigger, and the mask occupies
\RsevChanceOccPct{} per cent of the searched band, giving
\emph{\RsevChanceExp{} expected unattributed Class~A crossings against
\RsevChanceObsAUnattr{} observed}. The two are not competing readings of one
quantity but the expectations for two different chains, the rank clause
being a further one-in-$(\NCtrl+1)$ filter; the observed population is
accounted for on either.

The screen's rank bias runs against us and not for us: it places the star
\emph{high} in its own control distribution, so it inflates star-first
outcomes rather than suppressing them. Folding that measured factor in
raises the expectation to \PriMean{} and leaves the observed count the
expected one.

\emph{One trade is not available to us, and it should be stated plainly.} An
excess of this size needs only a small unmodelled tail in the control
ensemble, and Appendix~\ref{app:radius} shows that one exists. But the same
noise model sets the injection recovery fractions, so a tail large enough to
remove the excess would make the quoted completeness optimistic by roughly
the same factor. The two cannot be traded independently, and neither
correction has been applied.

\subsection{The control ensemble: its calibration, and its resolution}
%% (subsection label retired: nothing cites it)

The screen's geometry is set in \S\ref{sec:statistic}; its calibration and
its resolution are the subject here. Under that geometry the stellar ranks
are consistent with the uniform distribution exchangeability implies, in
both array strata --- \VerTwmAfterN{} 12\,m windows and \VerAcaAfterN{} ACA
windows, the latter with median rank \VerAcaAfterMed{} against the one half
expected --- over the \VerNTested{} of \NWindows{} windows that retain a
control vector, the remaining \VerNUnmatched{} having none. The ACA windows are screened on the
geometry measured from each observation's own antenna table and baselines
rather than on a 12\,m beam, and the array configuration does not enter the
sensitivity. Two cautions keep this from being a proof of exchangeability.
With samples of this size a shift of less than about \VerKsPowerAca{} in
the rank distribution could not have been seen at all, so these strata are
consistent with uniform rather than demonstrably uniform;
and the diagnostic most directly sensitive to contamination of the inner
controls, the excess of the innermost control bin over the outer annulus, is
\VerAcaAfterExc, with an interval of \VerAcaAfterExcLo{} to
\VerAcaAfterExcHi{} that spans zero and therefore neither confirms nor
excludes a residual effect.

\emph{Neither form is a significance test, and the reason is structural.}
Ranking the star against \NCtrl{} control positions drawn from the same field
is a powerful way to decide which windows deserve examination, but
$1/(\NCtrl+1)=\RankFloorVal$ is the \emph{resolution} of the ensemble --- the
smallest rank it can express --- and not the probability that noise produced
the window. Against a survey-wide Bonferroni scale
%
\begin{equation}
\alpha_{\rm survey}=\frac{0.05}{\NWindows}=\BonfScaleCalc,
\qquad
\frac{\RankFloorVal}{\alpha_{\rm survey}}=\BonfRankRatio,
\label{eq:bonf}
\end{equation}
%
that floor is \BonfRankRatio{} times too coarse, and correlation widens it:
neighbouring controls are correlated below the synthesised-beam scale, so a
window's annulus carries only $N_{\rm eff}\simeq30$--$43$ independent spatial
trials and the ensemble resolves no better than \RankEffLo--\RankEffHi.
Drawing more controls does not rescue it, the primary beam bounding how many
independent positions a single pointing contains. \textbf{The floor of
Eq.~\ref{eq:bonf} is the design's central limitation: no single window in this survey can be
significant on the rank alone, which is why a technosignature claim here
requires independent recurrence or independent astrophysical confirmation.}
What a visibility-domain statistic buys instead is measurable on this survey's
own data: in the image plane every rank-passing window sits at rank 1 of
$\NCtrl+1$ with no dynamic range left to separate them, whereas the
visibility fit separates the same windows by \VgSep$\sigma$ and its strongest
reaches \VgVisMax$\sigma$.

\subsection{The trials budget}

A window holds \CellsWinMin--\CellsWinMax{} channel$\times$drift cells
(median \CellsWinMed), the survey \CellsTotal{} in total, and they are not
independent trials: channel smoothing and a drift grid oversampled two to
one together over-count them by \FaOverCount. The control ensembles
calibrate that budget without a Gaussian assumption. The same arithmetic
predicts \NCtrlWinExp{} of the \NWindows{} windows carrying a control above
$T_{\rm trig}$ and \NCtrlWinObs{} do, agreeing to \CtrlWinAgreePct{} per
cent; corrected for the over-count, the observed control maxima sit
$+\FaExcessFine$ and $+\FaExcessCoarse$ per cent above prediction in the
fine and coarse classes independently. Non-Gaussianity of that size is why
the operative statement is the empirical rank and not the Gaussian
budget.

\emph{The residual on-star excess is not zero, and it is reported as an open
systematic rather than explained away.} Correlation cannot move the expected
\emph{number} above the trigger, only its clumping, so the excess over an
iid-Gaussian budget divides in two: the cells within one crossing window, a
factor \CellsPerCrossRatio{} that the frequency-axis over-count predicts
because one real excursion is counted about three times, and the number of
crossing windows, a factor \CrossWinRatio{} with a one-sided Poisson
probability of \CrossWinPoisson. The residual runs one way and carries the
same sign at the controls. The crossing list is dominated by
CO($2{\to}1$) at 230.5\,GHz
toward disc hosts, and a residual excess of chance crossings at the
tens-of-per-cent level cannot be excluded.

\emph{One stratum departs from uniformity and it is the one that matters.}
The fine-channel windows' ranks are not uniform, at $\pv{\KsFineP}$ over
$n=\NFine$ windows, because \NStageOneFine{} of them contain an on-star peak
above their own control maximum against \ExpStageOneFine{} expected; the coarse
stratum, three times larger and predicting \ExpStageOneCoarse{} such windows,
contains \NStageOneCoarse. A 15.6--31.25\,MHz channel dilutes a narrow
circumstellar line, so the class that can resolve one is the class that finds
it, and the same reasoning says a genuine unresolved carrier would appear in
Class~A first. The departure is nonetheless genuine and survives removal of
every fine window carrying a stellar crossing, so fine-window rank
probabilities are optimistic for the class as a whole and not only where a
crossing sits.

\emph{What is left is bounded, and the bounds are null.} Treating each
control as a pseudo-star and ranking it against the other \NCtrl-1 gives
\NPseudo{} ranks that are uniform to the resolution the ensemble can
express, and the real stellar ranks agree with them at $\pv{\StarPseudoP}$.
Any systematic radial suppression of the controls is below
\RingRhoBoundPct{} per cent in the mean, and the star's own residual
variance equals the controls' to better than a per cent in the rank-passing
windows. Scrambling each window so that nothing coherent survives along a
drift track carries \LNNWindows{} windows past the rank floor and changes no
disposition; it fails outright only on HD~48370 Band~6, whose extended
emission and imaging residuals are coherent across integrations and are
destroyed by the scramble itself, so no scrambled-noise probability is
quoted there. The remaining stratifications --- by per-control noise and by
gain --- cannot be made at all, because neither is kept in the retained
products.

%% ------------------------------------------------------------------
%% REQUEST (appendix-triage -> numbers):
%%   This appendix now leads with the TRIGGER + RANK expectation
%%   (\ClustStageMean against \ClustStageNObs) because that is the chain the
%%   revision adopts, and reports the trigger+localisation+attribution
%%   expectation (\RsevChanceExp against \RsevChanceObsAUnattr) second, as a
%%   different conditioning rather than a superseded one.  If the \Rsev
%%   family is being retired with the x2.88 sensitivity (integration queue
%%   item 3), say so and this paragraph goes; if it survives, it must not be
%%   reachable as "the" sensitivity.  One of the two has to be chosen, and it
%%   is not an appendix decision.
%% REQUEST (appendix-triage -> integrator):
%%   Floats deleted from this appendix and moved to the data release:
%%   tab:ctrlsep (control separations), tab:excess (the cell-excess
%%   decomposition), tab:arraysplit (the array split) and fig:ctrldiag
%%   (control diagnostics).  None is cited from the body.  tab_ctrlsep_v400,
%%   tab_falsealarm_v399 and the control_diagnostics figure are now
%%   uninputted; retire_macros will strip their macros.
%% ------------------------------------------------------------------
